\iftestbuild % TEST-ONLY-BEGIN
\setcounter{section}{1}\setcounter{theorem}{21}
\fi % TEST-ONLY-END
% Source: book pp. 197--210 (PDF 211--224); no solutions for this section.
\section{Orthonormal Bases}

\subsection{Orthonormal Lists and the Gram--Schmidt Procedure}

\begin{definition}[Orthonormal]\label{ax:6.22}\leavevmode
\begin{itemize}
\item A list of vectors is called \emph{orthonormal} if each vector in the list
  has norm $1$ and is orthogonal to all the other vectors in the list.
\item In other words, a list $e_1,\dots,e_m$ of vectors in $V$ is orthonormal if
  \[ \ip{e_j}{e_k} = \begin{cases} 1 & \text{if } j=k,\\
                                   0 & \text{if } j\ne k \end{cases} \]
  for all $j,k\in\{1,\dots,m\}$.
\end{itemize}
\end{definition}

\begin{example}[Orthonormal lists]\label{ax:6.23}\leavevmode
\begin{parts}
\item The standard basis of $\F^n$ is an orthonormal list.
\item $\bigl(1/\sqrt3,1/\sqrt3,1/\sqrt3\bigr),
  \bigl(-1/\sqrt2,1/\sqrt2,0\bigr)$ is an orthonormal list
  in $\F^3$.
\item $\bigl(1/\sqrt3,1/\sqrt3,1/\sqrt3\bigr),
  \bigl(-1/\sqrt2,1/\sqrt2,0\bigr),
  \bigl(1/\sqrt6,1/\sqrt6,-2/\sqrt6\bigr)$ is an
  orthonormal list in $\F^3$.
\item Suppose $n$ is a positive integer. Then, as Exercise~4 asks you to verify,
  \[ \frac{1}{\sqrt{2\pi}}, \frac{\cos x}{\sqrt\pi}, \frac{\cos 2x}{\sqrt\pi},
     \dots, \frac{\cos nx}{\sqrt\pi}, \frac{\sin x}{\sqrt\pi},
     \frac{\sin 2x}{\sqrt\pi}, \dots, \frac{\sin nx}{\sqrt\pi} \]
  is an orthonormal list of vectors in $C[-\pi,\pi]$, the vector space of
  continuous real-valued functions on $[-\pi,\pi]$ with inner product
  \[ \ip{f}{g} = \int_{-\pi}^{\pi} fg. \]
  The orthonormal list above is often used for modeling periodic phenomena,
  such as tides.
\item Suppose we make $\Poly_2(\R)$ into an inner product space using the inner
  product given by
  \[ \ip{p}{q} = \int_{-1}^{1} pq \]
  for all $p,q\in\Poly_2(\R)$. The standard basis $1,x,x^2$ of $\Poly_2(\R)$ is
  not an orthonormal list because the vectors in that list do not have norm $1$.
  Dividing each vector by its norm gives the list $1/\sqrt2$,
  $\sqrt{3/2}\,x$, $\sqrt{5/2}\,x^2$, in which each vector has norm $1$, and the
  second vector is orthogonal to the first and third vectors. However, the
  first and third vectors are not orthogonal. Thus this is not an orthonormal
  list. Soon we will see how to construct an orthonormal list from the standard
  basis $1,x,x^2$ (see Example~\ref{ax:6.34}).
\end{parts}
\end{example}

Orthonormal lists are particularly easy to work with, as illustrated by the
next result.

\begin{result}[Norm of an orthonormal linear combination]\label{ax:6.24}
Suppose $e_1,\dots,e_m$ is an orthonormal list of vectors in $V$. Then
\[ \norm{a_1e_1+\cdots+a_me_m}^2 = \abs{a_1}^2+\cdots+\abs{a_m}^2 \]
for all $a_1,\dots,a_m\in\F$.
\end{result}

\begin{proof}
Because each $e_k$ has norm $1$, this follows from repeated applications of the
Pythagorean theorem (\ref{ax:6.12}).
\end{proof}

The result above has the following important corollary.

\begin{result}[Orthonormal lists are linearly independent]\label{ax:6.25}
Every orthonormal list of vectors is linearly independent.
\end{result}

\begin{proof}
Suppose $e_1,\dots,e_m$ is an orthonormal list of vectors in $V$ and
$a_1,\dots,a_m\in\F$ are such that
\[ a_1e_1+\cdots+a_me_m = 0. \]
Then $\abs{a_1}^2+\cdots+\abs{a_m}^2=0$ (by \ref{ax:6.24}), which means that all
the $a_k$'s are $0$. Thus $e_1,\dots,e_m$ is linearly independent.
\end{proof}

Now we come to an important inequality.

\begin{result}[Bessel's inequality]\label{ax:6.26}
Suppose $e_1,\dots,e_m$ is an orthonormal list of vectors in $V$. If $v\in V$
then
\[ \abs{\ip{v}{e_1}}^2+\cdots+\abs{\ip{v}{e_m}}^2 \le \norm{v}^2. \]
\end{result}

\begin{proof}
Suppose $v\in V$. Then
\[ v = \underbrace{\ip{v}{e_1}e_1+\cdots+\ip{v}{e_m}e_m}_{u}
     + \underbrace{v-\ip{v}{e_1}e_1-\cdots-\ip{v}{e_m}e_m}_{w}. \]

Let $u$ and $w$ be defined as in the equation above. If $k\in\{1,\dots,m\}$, then
$\ip{w}{e_k}=\ip{v}{e_k}-\ip{v}{e_k}\ip{e_k}{e_k}=0$. This implies that
$\ip{w}{u}=0$. The Pythagorean theorem now implies that
\begin{align*}
\norm{v}^2 &= \norm{u}^2+\norm{w}^2\\
           &\ge \norm{u}^2\\
           &= \abs{\ip{v}{e_1}}^2+\cdots+\abs{\ip{v}{e_m}}^2,
\end{align*}
where the last line comes from \ref{ax:6.24}.
\end{proof}

The next definition introduces one of the most useful concepts in the study of
inner product spaces.

\begin{definition}[Orthonormal basis]\label{ax:6.27}
An \emph{orthonormal basis} of $V$ is an orthonormal list of vectors in $V$ that
is also a basis of $V$.
\end{definition}

For example, the standard basis is an orthonormal basis of $\F^n$.

\begin{result}[Orthonormal lists of the right length are orthonormal
bases]\label{ax:6.28}
Suppose $V$ is finite-dimensional. Then every orthonormal list of vectors in $V$
of length $\dim V$ is an orthonormal basis of $V$.
\end{result}

\begin{proof}
By \ref{ax:6.25}, every orthonormal list of vectors in $V$ is linearly
independent. Thus every such list of the right length is a basis---see
\ref{ax:2.38}.
\end{proof}

\begin{example}[An orthonormal basis of $\F^4$]\label{ax:6.29}
As mentioned above, the standard basis is an orthonormal basis of $\F^4$. We now
show that
\[ (1/2,1/2,1/2,1/2),
   (1/2,1/2,-1/2,-1/2),
   (1/2,-1/2,-1/2,1/2),
   (-1/2,1/2,-1/2,1/2) \]
is also an orthonormal basis of $\F^4$.

We have
\[ \bigl\lVert(1/2,1/2,1/2,1/2)\bigr\rVert
   = \sqrt{1/2^2+1/2^2+1/2^2+1/2^2} = 1. \]
Similarly, the other three vectors in the list above also have norm $1$.

Note that
\[ \bigl\langle(1/2,1/2,1/2,1/2),(1/2,1/2,-1/2,-1/2)\bigr\rangle
   = 1/2\cdot1/2+1/2\cdot1/2+1/2\cdot(-1/2)+1/2\cdot(-1/2) = 0. \]
Similarly, the inner product of any two distinct vectors in the list above also
equals $0$.

Thus the list above is orthonormal. Because we have an orthonormal list of
length four in the four-dimensional vector space $\F^4$, this list is an
orthonormal basis of $\F^4$ (by \ref{ax:6.28}).
\end{example}

In general, given a basis $e_1,\dots,e_n$ of $V$ and a vector $v\in V$, we know
that there is some choice of scalars $a_1,\dots,a_n\in\F$ such that
\[ v = a_1e_1+\cdots+a_ne_n. \]
Computing the numbers $a_1,\dots,a_n$ that satisfy the equation above can be a
long computation for an arbitrary basis of $V$. The next result shows, however,
that this is easy for an orthonormal basis---just take $a_k=\ip{v}{e_k}$.

\marginnote{The formula below for $\norm{v}$ is called Parseval's
identity. It was published in 1799 in the context of Fourier series.}%
Notice how the next result makes each inner product space of dimension $n$
behave like $\F^n$, with the role of the coordinates of a vector in $\F^n$
played by $\ip{v}{e_1},\dots,\ip{v}{e_n}$.

\begin{result}[Writing a vector as a linear combination of an orthonormal
basis]\label{ax:6.30}
Suppose $e_1,\dots,e_n$ is an orthonormal basis of $V$ and $u,v\in V$. Then
\begin{parts}
\item $v=\ip{v}{e_1}e_1+\cdots+\ip{v}{e_n}e_n$;
\item $\norm{v}^2=\abs{\ip{v}{e_1}}^2+\cdots+\abs{\ip{v}{e_n}}^2$;
\item $\ip{u}{v}=\ip{u}{e_1}\overline{\ip{v}{e_1}}+\cdots
  +\ip{u}{e_n}\overline{\ip{v}{e_n}}$.
\end{parts}
\end{result}

\begin{proof}
Because $e_1,\dots,e_n$ is a basis of $V$, there exist scalars $a_1,\dots,a_n$
such that
\[ v = a_1e_1+\cdots+a_ne_n. \]
Because $e_1,\dots,e_n$ is orthonormal, taking the inner product of both sides of
this equation with $e_k$ gives $\ip{v}{e_k}=a_k$. Thus (a) holds.

Now (b) follows immediately from (a) and \ref{ax:6.24}.

Take the inner product of $u$ with each side of (a) and then get (c) by using
conjugate linearity [\ref{ax:6.6}(d) and \ref{ax:6.6}(e)] in the second slot of
the inner product.
\end{proof}

\begin{example}[Finding coefficients for a linear combination]\label{ax:6.31}
Suppose we want to write the vector $(1,2,4,7)\in\F^4$ as a linear combination
of the orthonormal basis
\[ (1/2,1/2,1/2,1/2),
   (1/2,1/2,-1/2,-1/2),
   (1/2,-1/2,-1/2,1/2),
   (-1/2,1/2,-1/2,1/2) \]
of $\F^4$ from Example~\ref{ax:6.29}. Instead of solving a system of four linear
equations in four unknowns, as typically would be required if we were working
with a nonorthonormal basis, we simply evaluate four inner products and use
\ref{ax:6.30}(a), getting that $(1,2,4,7)$ equals
\[ 7(1/2,1/2,1/2,1/2)
   -4(1/2,1/2,-1/2,-1/2)
   +(1/2,-1/2,-1/2,1/2)
   +2(-1/2,1/2,-1/2,1/2). \]
\end{example}

Now that we understand the usefulness of orthonormal bases, how do we go about
finding them? For example, does $\Poly_m(\R)$ with inner product as in
\ref{ax:6.3}(c) have an orthonormal basis? The next result will lead to answers
to these questions.

\marginnote{J\o rgen Gram (1850--1916) and Erhard Schmidt (1876--1959)
popularized this algorithm that constructs orthonormal lists.}%
The algorithm used in the next proof is called the \emph{Gram--Schmidt
procedure}. It gives a method for turning a linearly independent list into an
orthonormal list with the same span as the original list.

\begin{result}[Gram--Schmidt procedure]\label{ax:6.32}
Suppose $v_1,\dots,v_m$ is a linearly independent list of vectors in $V$. Let
$f_1=v_1$. For $k=2,\dots,m$, define $f_k$ inductively by
\[ f_k = v_k-\frac{\ip{v_k}{f_1}}{\norm{f_1}^2}f_1-\cdots
         -\frac{\ip{v_k}{f_{k-1}}}{\norm{f_{k-1}}^2}f_{k-1}. \]
For each $k=1,\dots,m$, let $e_k=\dfrac{f_k}{\norm{f_k}}$. Then
$e_1,\dots,e_m$ is an orthonormal list of vectors in $V$ such that
\[ \spn(v_1,\dots,v_k) = \spn(e_1,\dots,e_k) \]
for each $k=1,\dots,m$.
\end{result}

\begin{proof}
We will show by induction on $k$ that the desired conclusion holds. To get
started with $k=1$, note that because $e_1=f_1/\norm{f_1}$, we have
$\norm{e_1}=1$; also, $\spn(v_1)=\spn(e_1)$ because $e_1$ is a nonzero multiple
of $v_1$.

Suppose $1<k\le m$ and the list $e_1,\dots,e_{k-1}$ generated by \ref{ax:6.32}
is an orthonormal list such that\refstepcounter{theorem}\label{ax:6.33}
\[ \spn(v_1,\dots,v_{k-1}) = \spn(e_1,\dots,e_{k-1}). \tag*{\thetheorem} \]
Because $v_1,\dots,v_m$ is linearly independent, we have
$v_k\notin\spn(v_1,\dots,v_{k-1})$. Thus
$v_k\notin\spn(e_1,\dots,e_{k-1})=\spn(f_1,\dots,f_{k-1})$, which implies that
$f_k\ne0$. Hence we are not dividing by $0$ in the definition of $e_k$ given in
\ref{ax:6.32}. Dividing a vector by its norm produces a new vector with
norm $1$; thus $\norm{e_k}=1$.

Let $j\in\{1,\dots,k-1\}$. Then
\begin{align*}
\ip{e_k}{e_j}
  &= \frac{1}{\norm{f_k}\,\norm{f_j}}\ip{f_k}{f_j}\\
  &= \frac{1}{\norm{f_k}\,\norm{f_j}}
     \biggl\langle v_k-\frac{\ip{v_k}{f_1}}{\norm{f_1}^2}f_1-\cdots
     -\frac{\ip{v_k}{f_{k-1}}}{\norm{f_{k-1}}^2}f_{k-1}, f_j\biggr\rangle\\
  &= \frac{1}{\norm{f_k}\,\norm{f_j}}(\ip{v_k}{f_j}-\ip{v_k}{f_j})\\
  &= 0.
\end{align*}
Thus $e_1,\dots,e_k$ is an orthonormal list.

From the definition of $e_k$ given in \ref{ax:6.32}, we see that
$v_k\in\spn(e_1,\dots,e_k)$. Combining this information with \ref{ax:6.33}
shows that
\[ \spn(v_1,\dots,v_k) \subseteq \spn(e_1,\dots,e_k). \]
Both lists above are linearly independent (the $v$'s by hypothesis, and the
$e$'s by orthonormality and \ref{ax:6.25}). Thus both subspaces above have
dimension $k$, and hence they are equal, completing the induction step and thus
completing the proof.
\end{proof}

\begin{example}[An orthonormal basis of $\Poly_2(\R)$]\label{ax:6.34}
Suppose we make $\Poly_2(\R)$ into an inner product space using the inner product
given by
\[ \ip{p}{q} = \int_{-1}^{1} pq \]
for all $p,q\in\Poly_2(\R)$. We know that $1,x,x^2$ is a basis of $\Poly_2(\R)$,
but it is not an orthonormal basis. We will find an orthonormal basis of
$\Poly_2(\R)$ by applying the Gram--Schmidt procedure with $v_1=1$, $v_2=x$,
and $v_3=x^2$.

To get started, take $f_1=v_1=1$. Thus $\norm{f_1}^2=\int_{-1}^{1}1=2$. Hence
the formula in \ref{ax:6.32} tells us that
\[ f_2 = v_2-\frac{\ip{v_2}{f_1}}{\norm{f_1}^2}f_1
       = x-\frac{\ip{x}{1}}{\norm{f_1}^2} = x, \]
where the last equality holds because $\ip{x}{1}=\int_{-1}^{1}t\,dt=0$.

The formula above for $f_2$ implies that
$\norm{f_2}^2=\int_{-1}^{1}t^2\,dt=2/3$. Now the formula in \ref{ax:6.32}
tells us that
\[ f_3 = v_3-\frac{\ip{v_3}{f_1}}{\norm{f_1}^2}f_1
            -\frac{\ip{v_3}{f_2}}{\norm{f_2}^2}f_2
       = x^2-1/2\ip{x^2}{1}-3/2\ip{x^2}{x}x = x^2-1/3. \]

The formula above for $f_3$ implies that
\[ \norm{f_3}^2 = \int_{-1}^{1}(t^2-1/3)^2\,dt
   = \int_{-1}^{1}(t^4-2/3t^2+1/9)\,dt = 8/45. \]

Now dividing each of $f_1,f_2,f_3$ by its norm gives us the orthonormal list
\[ \sqrt{1/2},\ \sqrt{3/2}\,x,\
   \sqrt{45/8}(x^2-1/3). \]
The orthonormal list above has length three, which is the dimension of
$\Poly_2(\R)$. Hence this orthonormal list is an orthonormal basis of
$\Poly_2(\R)$ [by \ref{ax:6.28}].
\end{example}

Now we can answer the question about the existence of orthonormal bases.

\begin{result}[Existence of orthonormal basis]\label{ax:6.35}
Every finite-dimensional inner product space has an orthonormal basis.
\end{result}

\begin{proof}
Suppose $V$ is finite-dimensional. Choose a basis of $V$. Apply the
Gram--Schmidt procedure (\ref{ax:6.32}) to it, producing an orthonormal list of
length $\dim V$. By \ref{ax:6.28}, this orthonormal list is an orthonormal basis
of $V$.
\end{proof}

Sometimes we need to know not only that an orthonormal basis exists, but also
that every orthonormal list can be extended to an orthonormal basis. In the next
corollary, the Gram--Schmidt procedure shows that such an extension is always
possible.

\begin{result}[Every orthonormal list extends to an orthonormal
basis]\label{ax:6.36}
Suppose $V$ is finite-dimensional. Then every orthonormal list of vectors in $V$
can be extended to an orthonormal basis of $V$.
\end{result}

\begin{proof}
Suppose $e_1,\dots,e_m$ is an orthonormal list of vectors in $V$. Then
$e_1,\dots,e_m$ is linearly independent (by \ref{ax:6.25}). Hence this list can
be extended to a basis $e_1,\dots,e_m,v_1,\dots,v_n$ of $V$ (see
\ref{ax:2.32}). Now apply the Gram--Schmidt procedure (\ref{ax:6.32}) to
$e_1,\dots,e_m,v_1,\dots,v_n$, producing an orthonormal list
\[ e_1,\dots,e_m,f_1,\dots,f_n; \]
here the formula given by the Gram--Schmidt procedure leaves the first $m$
vectors unchanged because they are already orthonormal. The list above is an
orthonormal basis of $V$ by \ref{ax:6.28}.
\end{proof}

Recall that a matrix is called upper triangular if it looks like this:
\[ \begin{pmatrix} * & & *\\ & \ddots & \\ 0 & & * \end{pmatrix}, \]
where the $0$ in the matrix above indicates that all entries below the diagonal
equal $0$, and asterisks are used to denote entries on and above the diagonal.

In the last chapter, we gave a necessary and sufficient condition for an
operator to have an upper-triangular matrix with respect to some basis (see
\ref{ax:5.44}). Now that we are dealing with inner product spaces, we would like
to know whether there exists an \emph{orthonormal} basis with respect to which
we have an upper-triangular matrix. The next result shows that the condition for
an operator to have an upper-triangular matrix with respect to some orthonormal
basis is the same as the condition to have an upper-triangular matrix with
respect to an arbitrary basis.

\begin{result}[Upper-triangular matrix with respect to some orthonormal
basis]\label{ax:6.37}
Suppose $V$ is finite-dimensional and $T\in\Lin(V)$. Then $T$ has an
upper-triangular matrix with respect to some orthonormal basis of $V$ if and
only if the minimal polynomial of $T$ equals $(z-\lambda_1)\cdots(z-\lambda_m)$
for some $\lambda_1,\dots,\lambda_m\in\F$.
\end{result}

\begin{proof}
Suppose $T$ has an upper-triangular matrix with respect to some basis
$v_1,\dots,v_n$ of $V$. Thus $\spn(v_1,\dots,v_k)$ is invariant under $T$ for
each $k=1,\dots,n$ (see \ref{ax:5.39}).

Apply the Gram--Schmidt procedure to $v_1,\dots,v_n$, producing an orthonormal
basis $e_1,\dots,e_n$ of $V$. Because
\[ \spn(e_1,\dots,e_k) = \spn(v_1,\dots,v_k) \]
for each $k$ (see \ref{ax:6.32}), we conclude that $\spn(e_1,\dots,e_k)$ is
invariant under $T$ for each $k=1,\dots,n$. Thus, by \ref{ax:5.39}, $T$ has an
upper-triangular matrix with respect to the orthonormal basis $e_1,\dots,e_n$.
Now use \ref{ax:5.44} to complete the proof.
\end{proof}

\marginnote{Issai Schur (1875--1941) published a proof of the next result in
1909.}%
For complex vector spaces, the next result is an important application of the
result above. See Exercise~20 for a version of Schur's theorem that applies
simultaneously to more than one operator.

\begin{result}[Schur's theorem]\label{ax:6.38}
Every operator on a finite-dimensional complex inner product space has an
upper-triangular matrix with respect to some orthonormal basis.
\end{result}

\begin{proof}
The desired result follows from the second version of the fundamental theorem of
algebra (\ref{ax:4.13}) and \ref{ax:6.37}.
\end{proof}

\subsection{Linear Functionals on Inner Product Spaces}

Because linear maps into the scalar field $\F$ play a special role, we defined a
special name for them and their vector space in Section~3F. Those definitions
are repeated below in case you skipped Section~3F.

\begin{definition}[Linear functional, dual space, $V'$]\label{ax:6.39}\leavevmode
\begin{itemize}
\item A \emph{linear functional} on $V$ is a linear map from $V$ to $\F$.
\item The \emph{dual space} of $V$, denoted by $V'$, is the vector space of all
  linear functionals on $V$. In other words, $V'=\Lin(V,\F)$.
\end{itemize}
\end{definition}

\begin{example}[Linear functional on $\F^3$]\label{ax:6.40}
The function $\varphi\colon\F^3\to\F$ defined by
\[ \varphi(z_1,z_2,z_3) = 2z_1-5z_2+z_3 \]
is a linear functional on $\F^3$. We could write this linear functional in the
form
\[ \varphi(z) = \ip{z}{w} \]
for every $z\in\F^3$, where $w=(2,-5,1)$.
\end{example}

\begin{example}[Linear functional on $\Poly_5(\R)$]\label{ax:6.41}
The function $\varphi\colon\Poly_5(\R)\to\R$ defined by
\[ \varphi(p) = \int_{-1}^{1} p(t)\bigl(\cos(\pi t)\bigr)\,dt \]
is a linear functional on $\Poly_5(\R)$.
\end{example}

\marginnote{The next result is named in honor of Frigyes Riesz (1880--1956), who
proved several theorems early in the twentieth century that look very much like
the result below.}%
If $v\in V$, then the map that sends $u$ to $\ip{u}{v}$ is a linear functional
on $V$. The next result states that every linear functional on $V$ is of this
form. For example, we can take $v=(2,-5,1)$ in Example~\ref{ax:6.40}.

Suppose we make the vector space $\Poly_5(\R)$ into an inner product space by
defining $\ip{p}{q}=\int_{-1}^{1}pq$. Let $\varphi$ be as in
Example~\ref{ax:6.41}. It is not obvious that there exists $q\in\Poly_5(\R)$
such that
\[ \int_{-1}^{1} p(t)\bigl(\cos(\pi t)\bigr)\,dt = \ip{p}{q} \]
for every $p\in\Poly_5(\R)$ [we cannot take $q(t)=\cos(\pi t)$ because that
choice of $q$ is not an element of $\Poly_5(\R)$]. The next result tells us the
somewhat surprising result that there indeed exists a polynomial
$q\in\Poly_5(\R)$ such that the equation above holds for all $p\in\Poly_5(\R)$.

\begin{result}[Riesz representation theorem]\label{ax:6.42}
Suppose $V$ is finite-dimensional and $\varphi$ is a linear functional on $V$.
Then there is a unique vector $v\in V$ such that
\[ \varphi(u) = \ip{u}{v} \]
for every $u\in V$.
\end{result}

\begin{proof}
First we show that there exists a vector $v\in V$ such that
$\varphi(u)=\ip{u}{v}$ for every $u\in V$. Let $e_1,\dots,e_n$ be an orthonormal
basis of $V$. Then
\begin{align*}
\varphi(u) &= \varphi(\ip{u}{e_1}e_1+\cdots+\ip{u}{e_n}e_n)\\
  &= \ip{u}{e_1}\varphi(e_1)+\cdots+\ip{u}{e_n}\varphi(e_n)\\
  &= \bigl\langle u, \overline{\varphi(e_1)}e_1+\cdots
     +\overline{\varphi(e_n)}e_n\bigr\rangle
\end{align*}
for every $u\in V$, where the first equality comes from \ref{ax:6.30}(a). Thus
setting\refstepcounter{theorem}\label{ax:6.43}
\[ v = \overline{\varphi(e_1)}e_1+\cdots+\overline{\varphi(e_n)}e_n,
   \tag*{\thetheorem} \]
we have $\varphi(u)=\ip{u}{v}$ for every $u\in V$, as desired.

Now we prove that only one vector $v\in V$ has the desired behavior. Suppose
$v_1,v_2\in V$ are such that
\[ \varphi(u) = \ip{u}{v_1} = \ip{u}{v_2} \]
for every $u\in V$. Then
\[ 0 = \ip{u}{v_1}-\ip{u}{v_2} = \ip{u}{v_1-v_2} \]
for every $u\in V$. Taking $u=v_1-v_2$ shows that $v_1-v_2=0$. Thus $v_1=v_2$,
completing the proof of the uniqueness part of the result.
\end{proof}

\begin{example}[Computation illustrating Riesz representation
theorem]\label{ax:6.44}
Suppose we want to find a polynomial $q\in\Poly_2(\R)$
such that\refstepcounter{theorem}\label{ax:6.45}
\[ \int_{-1}^{1} p(t)\bigl(\cos(\pi t)\bigr)\,dt = \int_{-1}^{1} pq
   \tag*{\thetheorem} \]
for every polynomial $p\in\Poly_2(\R)$. To do this, we make $\Poly_2(\R)$ into
an inner product space by defining $\ip{p}{q}$ to be the right side of the
equation above for $p,q\in\Poly_2(\R)$. Note that the left side of the equation
above does not equal the inner product in $\Poly_2(\R)$ of $p$ and the function
$t\mapsto\cos(\pi t)$ because this last function is not a polynomial.

Define a linear functional $\varphi$ on $\Poly_2(\R)$ by letting
\[ \varphi(p) = \int_{-1}^{1} p(t)\bigl(\cos(\pi t)\bigr)\,dt \]
for each $p\in\Poly_2(\R)$. Now use the orthonormal basis from
Example~\ref{ax:6.34} and apply formula \ref{ax:6.43} from the proof of the
Riesz representation theorem to see that if $p\in\Poly_2(\R)$, then
$\varphi(p)=\ip{p}{q}$, where
\begin{align*}
q(x) &= \biggl(\int_{-1}^{1}\sqrt{1/2}\cos(\pi t)\,dt\biggr)\sqrt{1/2}
        +\biggl(\int_{-1}^{1}\sqrt{3/2}\,t\cos(\pi t)\,dt\biggr)
        \sqrt{3/2}\,x\\
     &\quad+\biggl(\int_{-1}^{1}\sqrt{45/8}(t^2-1/3)
        \cos(\pi t)\,dt\biggr)\sqrt{45/8}(x^2-1/3).
\end{align*}
A bit of calculus applied to the equation above shows that
\[ q(x) = \tfrac{15}{2\pi^2}(1-3x^2). \]

The same procedure shows that if we want to find $q\in\Poly_5(\R)$ such that
\ref{ax:6.45} holds for all $p\in\Poly_5(\R)$, then we should take
\[ q(x) = \tfrac{105}{8\pi^4}\Bigl((27-2\pi^2)+(24\pi^2-270)x^2
          +(315-30\pi^2)x^4\Bigr). \]
\end{example}

Suppose $V$ is finite-dimensional and $\varphi$ a linear functional on $V$. Then
\ref{ax:6.43} gives a formula for the vector $v$ that satisfies
\[ \varphi(u) = \ip{u}{v} \]
for all $u\in V$. Specifically, we have
\[ v = \overline{\varphi(e_1)}e_1+\cdots+\overline{\varphi(e_n)}e_n. \]
The right side of the equation above seems to depend on the orthonormal basis
$e_1,\dots,e_n$ as well as on $\varphi$. However, \ref{ax:6.42} tells us that
$v$ is uniquely determined by $\varphi$. Thus the right side of the equation
above is the same regardless of which orthonormal basis $e_1,\dots,e_n$ of $V$
is chosen.

For two additional different proofs of the Riesz representation theorem, see
\ref{ax:6.58} and also Exercise~13 in Section~6C.

% ------------------------------------------------------------------ exercises
\exercisesheading{6B}

\begin{exercise}
Suppose $e_1,\dots,e_m$ is a list of vectors in $V$ such that
\[ \norm{a_1e_1+\cdots+a_me_m}^2 = \abs{a_1}^2+\cdots+\abs{a_m}^2 \]
for all $a_1,\dots,a_m\in\F$. Show that $e_1,\dots,e_m$ is an orthonormal list.
\exnote{This exercise provides a converse to \ref{ax:6.24}.}
\end{exercise}

\begin{exercise}\leavevmode
\begin{parts}
\item Suppose $\theta\in\R$. Show that both
  \[ (\cos\theta,\sin\theta), (-\sin\theta,\cos\theta) \quad\text{and}\quad
     (\cos\theta,\sin\theta), (\sin\theta,-\cos\theta) \]
  are orthonormal bases of $\R^2$.
\item Show that each orthonormal basis of $\R^2$ is of the form given by one of
  the two possibilities in (a).
\end{parts}
\end{exercise}

\begin{exercise}
Suppose $e_1,\dots,e_m$ is an orthonormal list in $V$ and $v\in V$. Prove that
\[ \norm{v}^2 = \abs{\ip{v}{e_1}}^2+\cdots+\abs{\ip{v}{e_m}}^2
   \iff v\in\spn(e_1,\dots,e_m). \]
\end{exercise}

\begin{exercise}
Suppose $n$ is a positive integer. Prove that
\[ \frac{1}{\sqrt{2\pi}}, \frac{\cos x}{\sqrt\pi}, \frac{\cos 2x}{\sqrt\pi},
   \dots, \frac{\cos nx}{\sqrt\pi}, \frac{\sin x}{\sqrt\pi},
   \frac{\sin 2x}{\sqrt\pi}, \dots, \frac{\sin nx}{\sqrt\pi} \]
is an orthonormal list of vectors in $C[-\pi,\pi]$, the vector space of
continuous real-valued functions on $[-\pi,\pi]$ with inner product
\[ \ip{f}{g} = \int_{-\pi}^{\pi} fg. \]
\exnote{Hint: The following formulas should help.
\begin{align*}
(\sin x)(\cos y) &= \frac{\sin(x-y)+\sin(x+y)}{2}\\
(\sin x)(\sin y) &= \frac{\cos(x-y)-\cos(x+y)}{2}\\
(\cos x)(\cos y) &= \frac{\cos(x-y)+\cos(x+y)}{2}
\end{align*}}
\end{exercise}

\begin{exercise}
Suppose $f\colon[-\pi,\pi]\to\R$ is continuous. For each nonnegative integer
$k$, define
\[ a_k = \tfrac{1}{\sqrt\pi}\int_{-\pi}^{\pi} f(x)\cos(kx)\,dx
   \quad\text{and}\quad
   b_k = \tfrac{1}{\sqrt\pi}\int_{-\pi}^{\pi} f(x)\sin(kx)\,dx. \]
Prove that
\[ \frac{a_0^2}{2}+\sum_{k=1}^{\infty}(a_k^2+b_k^2) \le \int_{-\pi}^{\pi} f^2. \]
\exnote{The inequality above is actually an equality for all continuous
functions $f\colon[-\pi,\pi]\to\R$. However, proving that this inequality is an
equality involves Fourier series techniques beyond the scope of this book.}
\end{exercise}

\begin{exercise}
Suppose $e_1,\dots,e_n$ is an orthonormal basis of $V$.
\begin{parts}
\item Prove that if $v_1,\dots,v_n$ are vectors in $V$ such that
  \[ \norm{e_k-v_k} < \frac{1}{\sqrt n} \]
  for each $k$, then $v_1,\dots,v_n$ is a basis of $V$.
\item Show that there exist $v_1,\dots,v_n\in V$ such that
  \[ \norm{e_k-v_k} \le \frac{1}{\sqrt n} \]
  for each $k$, but $v_1,\dots,v_n$ is not linearly independent.
\end{parts}
\exnote{This exercise states in \textup{(a)} that an appropriately small
perturbation of an orthonormal basis is a basis. Then \textup{(b)} shows that
the number $1/\sqrt n$ on the right side of the inequality in \textup{(a)}
cannot be higher.}
\end{exercise}

\begin{exercise}
Suppose $T\in\Lin(\R^3)$ has an upper-triangular matrix with respect to the
basis $(1,0,0)$, $(1,1,1)$, $(1,1,2)$. Find an orthonormal basis of $\R^3$ with
respect to which $T$ has an upper-triangular matrix.
\end{exercise}

\begin{exercise}
Make $\Poly_2(\R)$ into an inner product space by defining
$\ip{p}{q}=\int_0^1 pq$ for all $p,q\in\Poly_2(\R)$.
\begin{parts}
\item Apply the Gram--Schmidt procedure to the basis $1,x,x^2$ to produce an
  orthonormal basis of $\Poly_2(\R)$.
\item The differentiation operator (the operator that takes $p$ to $p'$) on
  $\Poly_2(\R)$ has an upper-triangular matrix with respect to the basis
  $1,x,x^2$, which is not an orthonormal basis. Find the matrix of the
  differentiation operator on $\Poly_2(\R)$ with respect to the orthonormal
  basis produced in (a) and verify that this matrix is upper triangular, as
  expected from the proof of \ref{ax:6.37}.
\end{parts}
\end{exercise}

\begin{exercise}
Suppose $e_1,\dots,e_m$ is the result of applying the Gram--Schmidt procedure to
a linearly independent list $v_1,\dots,v_m$ in $V$. Prove that
$\ip{v_k}{e_k}>0$ for each $k=1,\dots,m$.
\end{exercise}

\begin{exercise}
Suppose $v_1,\dots,v_m$ is a linearly independent list in $V$. Explain why the
orthonormal list produced by the formulas of the Gram--Schmidt procedure
(\ref{ax:6.32}) is the only orthonormal list $e_1,\dots,e_m$ in $V$ such that
$\ip{v_k}{e_k}>0$ and $\spn(v_1,\dots,v_k)=\spn(e_1,\dots,e_k)$ for each
$k=1,\dots,m$.
\exnote{The result in this exercise is used in the proof of \ref{ax:7.58}.}
\end{exercise}

\begin{exercise}
Find a polynomial $q\in\Poly_2(\R)$ such that
$p(1/2)=\int_0^1 pq$ for every $p\in\Poly_2(\R)$.
\end{exercise}

\begin{exercise}
Find a polynomial $q\in\Poly_2(\R)$ such that
\[ \int_0^1 p(x)\cos(\pi x)\,dx = \int_0^1 pq \]
for every $p\in\Poly_2(\R)$.
\end{exercise}

\begin{exercise}
Show that a list $v_1,\dots,v_m$ of vectors in $V$ is linearly dependent if and
only if the Gram--Schmidt formula in \ref{ax:6.32} produces $f_k=0$ for some
$k\in\{1,\dots,m\}$.
\exnote{This exercise gives an alternative to Gaussian elimination techniques
for determining whether a list of vectors in an inner product space is linearly
dependent.}
\end{exercise}

\begin{exercise}
Suppose $V$ is a real inner product space and $v_1,\dots,v_m$ is a linearly
independent list of vectors in $V$. Prove that there exist exactly $2^m$
orthonormal lists $e_1,\dots,e_m$ of vectors in $V$ such that
\[ \spn(v_1,\dots,v_k) = \spn(e_1,\dots,e_k) \]
for all $k\in\{1,\dots,m\}$.
\end{exercise}

\begin{exercise}
Suppose $\ip{\cdot}{\cdot}_1$ and $\ip{\cdot}{\cdot}_2$ are inner products on $V$
such that $\ip{u}{v}_1=0$ if and only if $\ip{u}{v}_2=0$. Prove that there is a
positive number $c$ such that $\ip{u}{v}_1=c\ip{u}{v}_2$ for every $u,v\in V$.
\exnote{This exercise shows that if two inner products have the same pairs of
orthogonal vectors, then each of the inner products is a scalar multiple of the
other inner product.}
\end{exercise}

\begin{exercise}
Suppose $V$ is finite-dimensional. Suppose $\ip{\cdot}{\cdot}_1$,
$\ip{\cdot}{\cdot}_2$ are inner products on $V$ with corresponding norms
$\norm{\cdot}_1$ and $\norm{\cdot}_2$. Prove that there exists a positive number
$c$ such that $\norm{v}_1\le c\norm{v}_2$ for every $v\in V$.
\end{exercise}

\begin{exercise}
Suppose $\F=\C$ and $V$ is finite-dimensional. Prove that if $T$ is an operator
on $V$ such that $1$ is the only eigenvalue of $T$ and $\norm{Tv}\le\norm{v}$
for all $v\in V$, then $T$ is the identity operator.
\end{exercise}

\begin{exercise}
Suppose $u_1,\dots,u_m$ is a linearly independent list in $V$. Show that there
exists $v\in V$ such that $\ip{u_k}{v}=1$ for all $k\in\{1,\dots,m\}$.
\end{exercise}

\begin{exercise}
Suppose $v_1,\dots,v_n$ is a basis of $V$. Prove that there exists a basis
$u_1,\dots,u_n$ of $V$ such that
\[ \ip{v_j}{u_k} = \begin{cases} 0 & \text{if } j\ne k,\\
                                 1 & \text{if } j=k. \end{cases} \]
\end{exercise}

\begin{exercise}
Suppose $\F=\C$, $V$ is finite-dimensional, and
$\mathcal{E}\subseteq\Lin(V)$ is such that
\[ ST = TS \]
for all $S,T\in\mathcal{E}$. Prove that there is an orthonormal basis of $V$
with respect to which every element of $\mathcal{E}$ has an upper-triangular
matrix.
\exnote{This exercise strengthens Exercise~9\textup{(b)} in Section~5E (in the
context of inner product spaces) by asserting that the basis in that exercise
can be chosen to be orthonormal.}
\end{exercise}

\begin{exercise}
Suppose $\F=\C$, $V$ is finite-dimensional, $T\in\Lin(V)$, and all eigenvalues
of $T$ have absolute value less than $1$. Let $\epsilon>0$. Prove that there
exists a positive integer $m$ such that $\norm{T^mv}\le\epsilon\norm{v}$ for
every $v\in V$.
\end{exercise}

\begin{exercise}
Suppose $C[-1,1]$ is the vector space of continuous real-valued functions on the
interval $[-1,1]$ with inner product given by
\[ \ip{f}{g} = \int_{-1}^{1} fg \]
for all $f,g\in C[-1,1]$. Let $\varphi$ be the linear functional on $C[-1,1]$
defined by $\varphi(f)=f(0)$. Show that there does not exist $g\in C[-1,1]$
such that
\[ \varphi(f) = \ip{f}{g} \]
for every $f\in C[-1,1]$.
\exnote{This exercise shows that the Riesz representation theorem
\textup{(\ref{ax:6.42})} does not hold on infinite-dimensional vector spaces
without additional hypotheses on $V$ and $\varphi$.}
\end{exercise}

\begin{exercise}
For all $u,v\in V$, define $d(u,v)=\norm{u-v}$.
\begin{parts}
\item Show that $d$ is a metric on $V$.
\item Show that if $V$ is finite-dimensional, then $d$ is a complete metric on
  $V$ (meaning that every Cauchy sequence converges).
\item Show that every finite-dimensional subspace of $V$ is a closed subset
  of $V$ (with respect to the metric $d$).
\end{parts}
\exnote{This exercise requires familiarity with metric spaces.}
\end{exercise}

\begin{remark*}[Orthogonality at the Supreme Court]\normalfont
Law professor Richard Friedman presenting a case before the U.S. Supreme Court
in 2010:
\begin{description}[font=\normalfont\itshape,style=sameline,leftmargin=2em,
  labelindent=0pt,itemsep=0pt,topsep=3pt]
\item[Mr. Friedman\textup{:}] I think that issue is entirely orthogonal to the
  issue here because the Commonwealth is acknowledging---
\item[Chief Justice Roberts\textup{:}] I'm sorry. Entirely what?
\item[Mr. Friedman\textup{:}] Orthogonal. Right angle. Unrelated. Irrelevant.
\item[Chief Justice Roberts\textup{:}] Oh.
\item[Justice Scalia\textup{:}] What was that adjective? I liked that.
\item[Mr. Friedman\textup{:}] Orthogonal.
\item[Chief Justice Roberts\textup{:}] Orthogonal.
\item[Mr. Friedman\textup{:}] Right, right.
\item[Justice Scalia\textup{:}] Orthogonal, ooh. (Laughter.)
\item[Justice Kennedy\textup{:}] I knew this case presented us a problem.
  (Laughter.)
\end{description}
\end{remark*}
